\chapter{Glossary and frequently asked questions}\label{app:glossary}

\section*{Glossary}
\begin{description}
  \item[Antagonism] share of interactions that are negative (FAULTANA).
  \item[Alignment] how well negative interactions follow the two-camp split (FAULTANA).
  \item[Beta distribution] distribution over a probability between 0 and 1; the Beta$(\alpha,\beta)$ prior updated by counts gives the posterior used for pair relations.
  \item[Bootstrap] resampling to estimate uncertainty (Section~\ref{sec:bootstrap}).
  \item[Censoring] the event of interest (leaving) is not observed within the window; right-censored: still active at the end; left-censored: began before the window.
  \item[CCDF] complementary cumulative distribution function, $P(X\ge x)$.
  \item[Balance, weak balance] a signed network is balanced if every triangle has three positive edges or one positive and two negative; it then splits into two camps. Weak balance also allows triangles with three negative edges and then several camps.
  \item[Frustration index] the least number of edges whose removal (or sign flip) makes a signed network balanced; hard to compute exactly.
  \item[Joint inference] a single test of whether the effect is zero in all specifications of a multiverse, by resampling under the null.
  \item[Popularity, novelty, root bias] the three ingredients of the thread-growth model of G\'omez et al.\ (2013).
  \item[Status theory] a positive link means the source regards the target as higher in a ranking; it explains directed signed links better than balance (Leskovec et al.\ 2010).
  \item[Cohort] group of users who first appeared in the same year.
  \item[Co-present null] random alternative target: author of a random earlier comment in the same article.
  \item[Confounder] a variable that affects both the exposure and the outcome and so distorts their comparison.
  \item[Cox model] regression for hazards; coefficients are log hazard ratios.
  \item[Dyad] an ordered pair of users $(A,B)$.
  \item[Edge] a reply from $A$ to $B$ (reply network) or a signed relation (stance network).
  \item[Faction] one of the two camps from FAULTANA.
  \item[Frustrated edge] an edge that violates balance: negative inside a camp or positive between camps.
  \item[Hazard ratio] ratio of instantaneous rates of leaving between two covariate values.
  \item[Immortal-time bias] bias from counting time during which the event could not have happened.
  \item[Kaplan--Meier] nonparametric survival estimate that handles censoring.
  \item[Landmark] a fixed time (day 30) from which follow-up starts, using only users still at risk.
  \item[Mantel--Haenszel] pooled odds ratio across strata.
  \item[Multiverse] the set of defensible analyses under alternative arbitrary choices.
  \item[Null model] a baseline that preserves some structure and randomises the thing being tested.
  \item[Odds ratio] ratio of odds; 1 means no difference.
  \item[Parquet] a compressed columnar file format.
  \item[Proxy] an observable variable standing in for an unobservable one (votes for stance).
  \item[Shrinkage] pulling a noisy estimate toward a population value.
  \item[Specification curve] plot of estimates across specifications, sorted by size.
  \item[Stratum] a subgroup within which a comparison is made.
  \item[Thread] a top-level comment and all the replies under it.
\end{description}

\section*{Frequently asked questions}
\begin{description}
  \item[Why not just use text?] The released data have none. Embeddings exist but are very large; the project works with votes.
  \item[Is a downvote really opposition?] Often, but not always: people downvote rudeness, off-topic or errors without disagreeing with the position. Treat votes as a signal of \emph{disapproval of the comment}; stance is an interpretation. The thesis should say so in its limitations.
  \item[Why 53\% vs 30\% but OR 4.43?] The odds ratio is computed within each comment and then pooled; the pooled shares are a coarser summary. See Section~\ref{sec:h1}.
  \item[Why is H1 the same as its positive-vote mirror?] Because for one table the two odds ratios are reciprocals.
  \item[Why 200 bootstrap replicates?] A balance of cost and stability for an interval; 1,000 would shift the endpoints very little.
  \item[Why landmark at 30 days?] A trade-off: later gives more exposure but removes more users; the multiverse tries 14 and 60 days.
  \item[Does HR 1.11 mean 11\% more users leave?] No: the rate of leaving is 11\% higher at any moment per SD of exposure.
  \item[Can I say that downvotes cause exit?] No. Observational, untimed votes, selected users.
  \item[What is the first thing to do before rerunning everything?] Commit the current state, so differences can be examined afterwards.
\end{description}
